\documentclass[twoside]{article}
\usepackage[utf8]{inputenc}

% The package supports other formats as well, but this template
% is primarily checked and updated for examinations at TUM BGDM
%\usepackage[uebung, answers]{tumbgdm}

\usepackage[uebung]{tumbgdm}
\usepackage{pythonhighlight}

\Professor{Prof. Dr. Martin Werner}
\Universitaet{Technical University of Munich}
\Institut{School of Engineering and Design}
\Lehrstuhl{Professorship Big Geospatial Data Management}
\Semester[WiTe 2024/25]{Winter Term 2024/25}
\Vorlesung{Principles of Programming}
\InstitutsWebsite{https://www.bgd.lrg.tum.de/}
\VorlesungsWebsite{https://www.moodle.tum.de/course/view.php?id=103704}

\nummerUebungsblatt{1}
\nummerErsteFrage{1}

\LearningOutcome{ Basic Python, Python Environments, Jupyter Notebooks, Code Editors, Command Line}

\begin{document}
\maketitle

\begin{task}{Environments and Jupyter Notebooks}{}{}

Create an environment and start a Jupyter Notebook:

\begin{lstlisting}[numbers=none]
python -m venv /path/to/my/envirnoment
source /path/to/my/envirnoment/bin/activate
pip install notebook
jupyter notebook
\end{lstlisting}

You can exit the environment by typing \textit{deactivate} in your shell.
If you want to use Anaconda, you may do this inside the "Anaconda Navigator" or execute the following lines:

\begin{lstlisting}[numbers=none]
conda create --name myfirstenv
conda activate myfirstenv
conda install notebook
jupyter notebook
\end{lstlisting}

You can exit the conda environment by typing \textit{conda deactivate} in your shell.

\end{task}

%Basic: https://pynative.com/python-basic-exercise-for-beginners/
\begin{task}{Basic Python Exercise}{}{}
Solve the following tasks inside a Jupyter Notebook. An example follows each task.
\begin{enumerate}
    \item print "Hello, World"
    \begin{lstlisting}[numbers=none]
    Output: Hello, World
    \end{lstlisting}
    
    \item multiplicate two numbers
    \begin{lstlisting}[numbers=none]
    Input: 53, 834
    Output: 44679
    \end{lstlisting}
    
    \item remove first n characters from a string
    \begin{lstlisting}[numbers=none]
    Input: "Hello", 2
    Output: llo
    \end{lstlisting}

    \clearpage

    \item check if the first and last number of a list are the same
    \begin{lstlisting}[numbers=none]
    Input: [5, 3, 1, 3, 4, 5]
    Output: True
    \end{lstlisting}
    
    
    \item return the count of a given substring from a given string
    \begin{lstlisting}[numbers=none]
    Input: "I have no idea what I am doing.", "I"
    Output: 2
    \end{lstlisting}
    
    \item check if a given number is a palindrome number
    \begin{lstlisting}[numbers=none]
    Input: 32123
    Output: True
    \end{lstlisting}
    
\end{enumerate}
\end{task}


\begin{task}{Python Scripts}{}{}

    Install a code editor (e.g., VS Code) and create a Python script (e.g., tutorial\_1.py). Implement Task 2 a. - f. with a main function and run it.
    The basic structure of your script may look like the following:
    
    \begin{lstlisting}[numbers=none]
        def main():
            # Your code here ...           
            return result

        if __name__=="__main__":
            main()
    \end{lstlisting}

\end{task}


\begin{task}{Command Line Python}{}{}

    Run your Python script from Task 3 in the command prompt (Windows) or terminal (Linux) by typing \textit{python /path/to/tutorial\_1.py}.
    Then, do Task 2 a. - b. without a script by running Python directly in the shell.
    You can start Python from the command line by typing \textit{python}.

\end{task}

\end{document}
